%% =====================================================================
%%  B3-T-07-01 -- what B3 contributes to "Posing as a Friend"
%%  -------------------------------------------------------------------
%%  LAYER A: APPEND-ONLY. Written once, then left alone. Material found
%%  only here sits in the `unique` slot and is protected by rule R10.
%% =====================================================================
%% @kind: source
%% @id: B3-T-07-01
%% @book: B3
%% @topic: T-07-01
%% @locator: Law 14, pp. 107--114
%% @status: distilled
%% @unique: yes
%% @integrated: yes
%% @updated: 2026-09-19

\dossier{B3}{T-07-01}{\bookshortB3}{Law 14, pp. 107--114}

\begin{distilled}
  \item Knowing about a rival is described as critical; spies are recommended for gathering information that keeps you a step ahead.
  \item Playing the spy yourself is recommended as better than employing one.
  \item Polite social encounters are the recommended venue, with indirect questions used to make people reveal weaknesses and intentions.
  \item There is no occasion that is not an opportunity for artful spying.
\end{distilled}

\begin{unique}
  \item The claim that every social occasion is usable, and the preference for collecting personally rather than through intermediaries, since a spy is an uncontrolled channel.
\end{unique}

\begin{terms}
  \item \emph{artful spying} --- information collection conducted inside ordinary social contact.
\end{terms}
